\subsection{Automated Hyperparameter Tuning}
\label{sec:automl_tuning}

Hyperparameter tuning is tedious. With sklearn, you manually define grid points
for each parameter---calculate appropriate values, balance resolution against
compute cost, manage complex dictionaries. For parameters like learning rates
that benefit from log-scale spacing, you need domain knowledge to pick good grid
points. MLPY automates this through declarative parameter spaces.

\subsubsection{Problem with Manual Grid Search}

Sklearn requires explicit grid points. Want to tune a Random Forest with
\texttt{n\_estimators} from 50 to 200? You manually write
\texttt{[50, 87, 125, 162, 200]}. For every parameter. For every tuning run.

\subsubsection{MLPY's Approach}

MLPY uses parameter types: \texttt{ParamInt} and \texttt{ParamFloat} for numeric
ranges, \texttt{ParamCategorical} for discrete choices, \texttt{ParamBool} for
flags. You specify ranges, not grid points. The tuner generates the grid based
on resolution.

\begin{lstlisting}[language=Python, caption={MLPY automated hyperparameter tuning for Random Forest}]
from mlpy.tasks import TaskClassif
from mlpy.learners.sklearn import LearnerRandomForestClassifier
from mlpy.measures import MeasureClassifAccuracy
from mlpy.resamplings import ResamplingCV
from mlpy.automl.tuning import (
    TunerGrid, ParamSet, ParamInt, ParamCategorical
)

# Create task
task = TaskClassif(data=df, target='species')

# Define learner
learner = LearnerRandomForestClassifier(random_state=42)

# Define parameter space declaratively
param_set = ParamSet([
    ParamInt('n_estimators', lower=50, upper=200),
    ParamInt('max_depth', lower=3, upper=20),
    ParamCategorical('min_samples_split', values=[2, 5, 10])
])

# Create tuner with desired resolution
tuner = TunerGrid(resolution=5)

# Run tuning
result = tuner.tune(
    learner=learner,
    task=task,
    resampling=ResamplingCV(folds=5),
    measure=MeasureClassifAccuracy(),
    param_set=param_set
)

# Access results
print(f"Best accuracy: {result.best_score:.4f}")
print(f"Best config: {result.best_config}")
\end{lstlisting}

The \texttt{resolution} parameter controls grid density: \texttt{resolution=5} generates 5 evenly-spaced values for each continuous parameter. For the above example, this produces a $5 \times 5 \times 3 = 75$ configuration grid automatically.

\subsubsection{Comparison with scikit-learn}

Table~\ref{tab:automl_code_comparison} compares MLPY and sklearn for the same
tuning task.

\begin{table}[ht]
\centering
\caption{Code comparison: MLPY vs scikit-learn for hyperparameter tuning}
\label{tab:automl_code_comparison}
\begin{tabular}{lcc}
\toprule
\textbf{Aspect} & \textbf{MLPY} & \textbf{scikit-learn} \\
\midrule
Lines of code & 12 & 14 \\
Parameter definition & Declarative ranges & Manual grid points \\
Grid generation & Automatic & Manual calculation \\
Task integration & Unified (Task + Learner) & Separate (X, y arrays) \\
Result object & Rich (\texttt{TuneResult}) & Basic (\texttt{GridSearchCV}) \\
\bottomrule
\end{tabular}
\end{table}

The key difference is parameter definition. Sklearn requires manual calculation:

\begin{lstlisting}[language=Python]
# scikit-learn: Manual grid point calculation
param_grid = {
    'n_estimators': [50, 87, 125, 162, 200],  # Manual
    'max_depth': [3, 7, 11, 15, 20],          # Manual
    'min_samples_split': [2, 5, 10]
}
\end{lstlisting}

MLPY's \texttt{ParamInt} and \texttt{ParamFloat} generate grids automatically,
cutting parameter definition code by about 40\%.

\subsubsection{Experimental Evaluation}

We compared MLPY and sklearn on Iris~\citep{fisher1936use}, tuning a Random
Forest over 18 configurations with 5-fold CV. Table~\ref{tab:automl_results}
shows comparable performance---similar time and accuracy.

\begin{table}[ht]
\centering
\caption{Hyperparameter tuning performance comparison on Iris dataset (150 samples, 4 features, 3 classes)}
\label{tab:automl_results}
\begin{tabular}{lccc}
\toprule
\textbf{Method} & \textbf{Time (s)} & \textbf{Best Accuracy} & \textbf{Configs} \\
\midrule
MLPY & 21.3 & 0.9600 & 18 \\
scikit-learn & 21.4 & 0.9667 & 18 \\
\bottomrule
\end{tabular}
\end{table}

\subsubsection{Advanced Features}

MLPY supports multiple tuning strategies (\texttt{TunerGrid}, \texttt{TunerRandom}),
log-scale parameters for learning rates, and callback integration for progress
tracking. The \texttt{TuneResult} object stores complete tuning history and
exports to DataFrames for analysis.

\subsubsection{Discussion}

The main advantage is eliminating boilerplate. Instead of manually calculating
grid points for every tuning run, you define parameter ranges once. This is
especially helpful for log-scale parameters and high-dimensional searches where
manual grids become unwieldy.

Performance matches sklearn, but the declarative API is cleaner and integrates
with MLPY's Task-Learner-Measure design.
